\scanpage{73}
\mspage{22}

\[d:X\longrightarrow X\times X=Y\] qui est donc dans $W$. \deleted{\unsure{Donc}} Mais \deleted{\unsure{alors}} $\pi_0(d)$ s'identifie à l'application diagonale \deleted{\unsure{pour}} $\pi_0(X)\longrightarrow\pi_0(X)\times\pi_0(X)$, qui n'est pas surjective puisque $\mathrm{card}\,\pi_0(X)\geq1$. Donc par la première partie, $W=\mathrm{Fl}(\mathrm{Cat})$.

Donc sous forme plus positive :

\uline{Corollaire} Soit $W\subset\mathrm{Fl}(\mathrm{Cat})$ qui satisfait à $W(1,2,3,4\,\mathrm b'')$, et telle que $W\ne\mathrm{Fl}(\mathrm{Cat})$. Alors $W$ \unsure{satisfais|satisfait} l'axiome de connexité $W(5)$, et par suite (th. 1, page 12) $W\subset W_\infty$.

III \unsure{Autres remarques|Autre remarque}. \unsure{Nom} des foncteurs liés à $W$.

Finalement, j'appellerai ``\uline{localiseur fondamental}'' une partie $W$ de $\mathrm{Fl}(\mathrm{Cat})$ satisfaisant les trois conditions suivantes :

(i) Saturation $W(1)$ (cf rappel p. 1)

(ii) Axiome de l'élément final $W(3)$ (page 3)

(iii) La variante affaiblie \inserted{$W(2\,\mathrm{bis})$} de \deleted{$W(2$} l'axiome de localisation $W(2)$ donnée dans le Cor. page 1 (anciennement c'était $\mathrm L\,4'$ -- dans P.S.)

Quand j'aurai besoin de plus sur $W$, je le dirai explicitement. Ces trois axiomes ont suffi pour \deleted{\unsure{que}} pratiquement tout ce que j'ai fait avec $W$ dans P.S., la seule exception étant un usage (si je me rappelle) unique, quelque part, d'une variante affaiblie de $W(2)$ (pour les foncteurs \uline{cartésiens} entre $S$-catégories \uline{fibrées} -- si \deleted{il induit} \unsure{dans} \unsure{toute} fibre il est $\in W$, il est lui-$=$ $\in W$).

On supposera donc par la suite donné un localiseur fondamental \deleted{$\unsure{W}$} \[W\subset\mathrm{Fl}(\mathrm{Cat}).\]

\uline{Proposition 1} Soit $f:X\longrightarrow Y$ \deleted{dans} \unsure{de} $\mathrm{Cat}$. Sont équiv. :

a) Pour tout \unsure{isom|isom.} local $Y'\longrightarrow Y$, $f':X'\longrightarrow Y'$ est $\in W$.

b) Même condition, dans le cas particulier $Y'=Y_{/y}\longrightarrow Y$, donc \deleted{$\unsure{\forall}$} $y\in Y$, $\underset{\unsure{\simeq}\,X\times_Y(Y_{/y})}{X_{/y}}$ est $W$-asphérique.

% archive sheet number: 73
